\chapter{汽轮机设备}

\begin{problem}
何谓汽轮机的级？级有哪几种类型？级是如何分类的？
\end{problem}

\begin{solution}
\begin{enumerate}
  \item \textbf{何谓汽轮机的级}："级"是汽轮机中最基本的工作单元。在结构上它是由静叶（喷嘴）和对应的动叶所组成；一列固定的喷嘴和与它配合的动叶片构成了汽轮机的基本作功单元，称为汽轮机的"级"。
  \item \textbf{级的分类及类型}：
  \begin{enumerate}
    \item 按照工作原理，汽轮机的级可分为冲动级和反动级两大类。其中冲动级又细分为纯冲动级和带反击度的冲动级。
    \item 按照工作特点，汽轮机的级还可分为速度级（包含双列复速级）和压力级。
  \end{enumerate}
\end{enumerate}
\end{solution}

\begin{problem}
汽轮机是如何分类的？大型机组中一般采用何种类型的汽轮机？汽轮机型号的各组成符号表明的含义是什么？
\end{problem}

\begin{solution}
\begin{enumerate}
  \item \textbf{汽轮机的分类}：
  \begin{enumerate}
    \item 按工作原理分：冲动式汽轮机、反动式汽轮机。
    \item 按热力特性分：凝汽式、背压式、调整抽汽式、抽汽背压式、中间再热式汽轮机。
    \item 按主蒸汽压力分：低压、中压、高压、超高压、亚临界压力、超临界压力、超超临界压力汽轮机。
    \item 此外还可按蒸汽流动方向（轴流/辐流式）、按汽缸数目（单缸/双缸/多缸）、按用途（电站/工业/船用）以及按功率大小进行分类。
  \end{enumerate}
  \item \textbf{大型机组采用的类型}：现代发电厂的大型机组一般采用多级汽轮机（如多缸、多排汽口、双轴系或高中压分流合缸结构），并广泛采用超临界或超超临界压力的中间再热凝汽式汽轮机。
  \item \textbf{汽轮机型号符号的含义}：国产汽轮机型号通常由四部分组成，一般表示形式为"汽轮机型式代号—额定功率—蒸汽参数—类型设计次序"。其中代号（如N代表凝汽式、B代表背压式）表明汽轮机型式；参数部分包含主蒸汽压力、温度及中间再热温度等数值。
\end{enumerate}
\end{solution}

\begin{problem}
汽轮机设备是由哪些设备构成的？各个设备的作用是什么？
\end{problem}

\begin{solution}
\begin{enumerate}
  \item \textbf{设备构成}：汽轮机设备主要包括汽轮机本体、调节保安及供油系统和辅助设备等。
  \item \textbf{各设备的作用}：
  \begin{enumerate}
    \item \textbf{汽轮机本体}：以蒸汽为工质，将蒸汽的热能转变为机械能，为发电机发电提供动力。
    \item \textbf{调节保安及供油系统}：控制汽轮机的进汽量，保障机组运行安全，并为系统提供润滑油和液压动力。
    \item \textbf{辅助设备}：包括凝汽设备、回热加热设备、水泵等。主要作用是建立排汽口真空、回收凝结水、对锅炉给水进行回热加热并除氧，完成热力循环。
  \end{enumerate}
\end{enumerate}
\end{solution}

\begin{problem}
汽轮机本体主要由哪些零部件组成？各零部件的结构特点及作用是什么？
\end{problem}

\begin{solution}
汽轮机本体包括固定部分（静子）和转动部分（转子）两大部分。

\textbf{静止部分（静子）}：
\begin{enumerate}
  \item \textbf{汽缸}：作为外壳将高温高压的蒸汽与大气隔开；包容隔板及转子等构成通流部分；承受零部件重量、内外压差及热应力；并加工有抽汽口用于回热。
  \item \textbf{隔板与喷嘴}：喷嘴（静叶片）相邻流道的作用是将蒸汽热能变为动能；隔板则用于固定喷嘴叶片并将汽缸内部空间分隔成若干个汽室。
  \item \textbf{汽封}：多采用梳齿形迷宫式结构。端部轴封防止蒸汽外漏或空气漏入真空区，通流部分汽封用于减少各级间的漏汽。
  \item \textbf{轴承}：采用滑动轴承。径向（支持）轴承用于支持转子重量并确定径向位置；推力轴承承担未平衡的轴向推力并确定转子的轴向位置。
\end{enumerate}

\textbf{转动部分（转子）}：
\begin{enumerate}
  \item \textbf{转子（主轴）}：现代大功率机组多采用整锻或焊接转子，用于汇集动叶栅上得到的机械能并传递扭矩。
  \item \textbf{叶轮}：装于主轴上，用来装置动叶片并将汽流扭矩传递给主轴。
  \item \textbf{动叶片}：由叶根、工作叶型和叶顶组成，作用是把蒸汽的动能转变成机械能，使转子旋转。
  \item \textbf{联轴器与盘车装置}：联轴器连接轴系并传递扭矩；盘车装置在停机或启动前驱动转子慢速旋转，防止转子受热不均产生热弯曲。
\end{enumerate}
\end{solution}

\begin{problem}
汽轮机级内能量转换过程分哪两步？分别在什么部件中完成？各步所转换能量的定量关系式是怎样的？
\end{problem}

\begin{solution}
\begin{enumerate}
  \item \textbf{第一步}：蒸汽的\textbf{热能转化为汽流的动能}。这一步在固定不动的喷嘴（静叶）中完成。
  \begin{enumerate}
    \item 定量关系式：根据稳定绝热流动的能量方程式，理想流速的定量关系为 \((h_0 - h_{1t}) \times 10^3 = \frac{1}{2}(c_{1t}^2 - c_0^2)\)；实际流速关系式为 \(c_1 = \varphi c_{1t} = 44.72 \varphi \sqrt{h_0 - h_{1t}}\)。
  \end{enumerate}
  \item \textbf{第二步}：汽流的\textbf{动能转变为轴的机械能}。这一步在高速旋转的\textbf{动叶流道}中完成。
  \begin{enumerate}
    \item 定量关系式：汽流作用于动叶上的轮周力 \(F_u = G(c_1 \cos\alpha_1 + c_2 \cos\alpha_2)\)；动叶损失的能量关系为 \(\delta h_b = \frac{1}{2}(w_1^2 - w_2^2) \times 10^{-3} = \frac{1}{2000}w_1^2(1 - \psi^2)\)。
  \end{enumerate}
\end{enumerate}
\end{solution}

\begin{problem}
级内存在哪些能量损失？它们是如何形成的？级内损失与级的理想焓降及有效烙降三者的关系如何？级效率的意义是什么？
\end{problem}

\begin{solution}
\begin{enumerate}
  \item \textbf{级内能量损失及其形成原因}：
  \begin{enumerate}
    \item \textbf{叶栅损失}：包括喷嘴和动叶的叶型摩擦损失、叶端（顶部和根部二次流）损失以及超音速产生的冲波损失。
    \item \textbf{余速损失}：蒸汽离开动叶栅时带走的排汽余速动能。
    \item \textbf{叶轮摩擦损失}：蒸汽粘性造成微团间及与叶轮的摩擦，以及叶轮间隙旋转流速不同形成的旋涡损失。
    \item \textbf{漏汽损失}：隔板、动叶顶部等动静间隙存在压差导致蒸汽漏走未作功。
    \item \textbf{其他损失}：还有扇形损失（叶栅沿高度节距/速度变化引起）、叶高损失、部分进汽损失以及湿汽损失（水珠消耗动能、阻止旋转）。
  \end{enumerate}
  \item \textbf{三者的关系}：级的有效焓降 \(\Delta H_i\) 等于级的理想焓降 \(\Delta H_t\) 减去所有的级内内部损失。
  \item \textbf{级效率的意义}：级的相对内效率表示级内能量转换的完善程度，是用来衡量汽轮机该级经济性的一个重要指标。
\end{enumerate}
\end{solution}

\begin{problem}
何谓汽轮机的内部损失？试述汽轮机的相对内效率所表明的意义？如何计算汽轮机的内功率？
\end{problem}

\begin{solution}
\begin{enumerate}
  \item \textbf{内部损失}：在汽轮机通流部分中，与蒸汽流动、能量转换有直接联系的能量损失统称为内部损失。它主要包括各级内的损失、进汽机构中的节流损失以及排汽管中的压力损失。
  \item \textbf{相对内效率的意义}：汽轮机的相对内效率（\(\eta_{0i}\)）是实际变为有用功的有效焓降 \(\Delta H_i\) 与总理想焓降 \(\Delta H_t\) 的比值。它表明了整个汽轮机内部能量转换的完善程度，反映了机组内部做功的经济性。
  \item \textbf{内功率的计算}：汽轮机的内功率 \(P_i\) 正比于蒸汽质量流量 \(D_0\) 与实际有效焓降 \(\Delta H_i\) 的乘积。计算公式为：
  \begin{equation*}
    P_i = \frac{D_0 \Delta H_i}{3600} = \frac{D_0 \Delta H_t \eta_{0i}}{3600}
  \end{equation*}
  （单位为 \unit{kW}）
\end{enumerate}
\end{solution}

\begin{problem}
汽轮机有哪些辅助设备？各辅助设备的作用、类型、结构特点及工作原理是什么？
\end{problem}

\begin{solution}
\textbf{凝汽设备（包含凝汽器、抽气器、冷却塔等）}：
\begin{enumerate}
  \item \textbf{凝汽器}：作用是建立并保持汽轮机排汽口的高度真空，并回收凝结水。多采用表面式，由外壳、水室、管板、冷却管束等组成。工作原理是由于蒸汽与水比体积相差极大，排汽在冷却管表面遇冷凝结成水，密闭空间体积骤缩从而形成高真空。
  \item \textbf{抽气器}：作用是不断抽出漏入的空气和不凝结气体以维持真空。分为射汽抽气器、射水抽气器和机械式真空泵（如水环式，结构简单启动快）。
\end{enumerate}

\textbf{回热加热设备（高压、低压加热器）}：
\begin{enumerate}
  \item \textbf{作用与类型}：利用汽轮机抽汽对锅炉给水进行加热，提高循环热效率。除除氧器外一般多采用表面式加热器，大型机组多为卧式结构。
  \item \textbf{结构特点}：主要由水室、U形管束（铜管或不锈钢管）、管板、隔板和壳体构成。高压加热器为了充分利用抽汽过热度，还设有蒸汽冷却段和疏水冷却段。
\end{enumerate}

\textbf{除氧器}：
\begin{enumerate}
  \item \textbf{作用与原理}：它是一级混合式加热器。作用是除去给水中的氧气及其他不凝结气体，防止热力设备腐蚀，同时汇集高压疏水等加以利用。
\end{enumerate}

\textbf{辅助水泵}：
\begin{enumerate}
  \item \textbf{凝结水泵}：从凝汽器热井中吸取凝结水输送至除氧器，工作在高度真空下，需采用抗汽蚀措施。
  \item \textbf{给水泵}：连续不断向锅炉输送高压给水，为离心式，多通过改变转速（小汽轮机驱动或液力耦合器）来控制流量。
  \item \textbf{循环水泵}：向凝汽器等供应冷却水，特点是水量大、扬程小，多采用轴流式或双吸离心式水泵。
\end{enumerate}
\end{solution}
